\subsection{Explicit hard sectors in the repo's counterexamples}
\label{ian:hard-sector-instances}

Status: proved for the listed instances, including covariance bounds with access to all physical spins and a domination repair for the soft-code law.

Throughout, \(m\ge2\) is even, \(n=2^m\), and
\[
 c_a(z)=(-1)^{a\cdot z},\qquad
 h(a)=\sum_{j=1}^{m/2}a_{2j-1}a_{2j}\pmod2,
 \qquad r=\frac{1-n^{-1/2}}2.
\]
The acceptance function is
\[
 P(x)=\sum_{a:h(a)=1}\widehat x(a)^2,\qquad
 \widehat x(a)=\frac1n\sum_zx_zc_a(z).
\]
It belongs to \([0,1]\) on the entire cube and is the acceptance
probability of one quantum query, by
Sections~\ref{qm:a-globally-bounded-quadratic-witness} and
\ref{qm:one-quantum-query-versus-logarithmically-many-classical-queries}.
In gadget models \(P\) uses only the visible spins, while classical
tests may query every physical spin.

\begin{lemma}[Covariance from the character-orbit calculation]
\label{ian:hard-sector-orbit-covariance}
Suppose a law admits a representation
\[
 X=\sigma c_A\eta,
\]
where \(A\) is uniform on \(\mathbb F_2^m\) and independent of side
information \(W\), and \(\sigma,\eta\) are determined by \(W\).
Suppose further that, after \(W\) is fixed, every depth-\(t\)
physical-spin test can be simulated using at most \(kt\) linear
queries to \(A\). Then
\[
 \E P=r,\qquad
 |\Cov(P,F)|\le
 c_{m,kt}:=\frac{2^{kt}-1}{2\sqrt n}
 \quad(F\in\mathcal F_t).
\]
If \(\eta\) always has fewer than \(n/32\) negative coordinates, then
\[
 \Var P\ge\left(\frac{97}{128}\right)^2r(1-r).
\]
Consequently, if in addition \(\eta\) always has fewer than \(n/32\) negative coordinates, then for \(n\ge16\) and \(kt\le m/2-2\),
\[
 \operatorname{Hard}\left(
 P,\mathcal L(\text{all physical spins});
 t,\frac{17}{128},\frac18\right).
\]
\end{lemma}

\begin{proof}
Fix \(W\), and put \(p_b=\widehat\eta(b)^2\). Parseval gives
\(\sum_bp_b=1\), and a change of Fourier index gives the exact formula
\[
 P(\sigma c_A\eta)=\sum_bp_bh(A+b).
\]
Thus its mean over \(A\) is \(r\). For every nonzero Walsh frequency
\(s\), the Fourier coefficient of \(h\) has magnitude
\(1/(2\sqrt n)\), by the bent-function calculation in
Section~\ref{qm:one-quantum-query-versus-logarithmically-many-classical-queries}.
Convolution by \((p_b)\) multiplies this coefficient by
\(\sum_bp_b(-1)^{s\cdot b}\), whose absolute value is at most one.

A leaf of a depth-\(kt\) linear-query tree conditions \(A\) onto an
affine subspace of codimension at most \(kt\). Its conditional mean
therefore satisfies
\[
 \left|\E[P-r\mid\text{leaf},W]\right|
 \le\frac{2^{kt}-1}{2\sqrt n}.
\]
Multiply by the leaf output in \([-1,1]\), sum over leaves, and average
over \(W\).

For the variance assertion,
\[
 p_0=\widehat\eta(0)^2\ge(15/16)^2=225/256.
\]
Write, as a function of uniform \(A\),
\[
 P=p_0h(A)+(1-p_0)R(A),
\]
where \(0\le R\le1\) and \(\E R=r\). Hence
\(\Var R\le r(1-r)\), and the reverse triangle inequality in \(L^2\)
gives
\[
 \sqrt{\Var(P\mid W)}
 \ge(2p_0-1)\sqrt{r(1-r)}
 \ge\frac{97}{128}\sqrt{r(1-r)}.
\]
The conditional mean is always \(r\), so averaging proves the bound.
For \(n\ge16\),
\[
 \left(\frac{97}{128}\right)^2r(1-r)
 \ge\frac{141135}{1048576}>\frac{17}{128},
 \qquad c_{m,kt}\le\frac18.
\]
Finally \(1/8<17/128\).
\end{proof}

This lemma obtains the covariance condition directly from the proof
of Lemma~\ref{res:orbit-lemma} in
\path{obstructions/local_codes/01_character_orbit_lemma.tex}.
An \(L^1\) approximation lower bound alone would not justify that
condition.

\begin{proposition}[The character code and arbitrary mixture weights]
\label{ian:hard-sector-character-instance}
Let \(n\ge16\) and \(t\le m/2-2\). Then
\[
 \operatorname{Hard}\left(P,U_{C_n};t,\frac3{16},\frac18\right).
\]
For the character-code Hamiltonian
\[
 E_n(x)=\Delta_n\mathbf1_{\{x\notin C_n\}},\qquad
 \Delta_n=\log(1+2^n/n),
\]
the set \(H=C_n\) has exact weight
\[
 p_{\beta,n}
 =\frac{n}{n+(2^n-n)e^{-\beta\Delta_n}}
 =w_{\beta,n}+(1-w_{\beta,n})\delta_n,
 \qquad \delta_n=n2^{-n},
\]
where
\[
 w_{\beta,n}
 =\frac{\delta_n(e^{\beta\Delta_n}-1)}
        {1+\delta_n(e^{\beta\Delta_n}-1)}.
\]
For every \(\beta>0\),
\[
 F_H-F
 =\frac1\beta\log\left(
 1+\frac{2^n-n}{n}e^{-\beta\Delta_n}\right).
\]
At \(\beta\ge1\), \(p_{\beta,n}\ge1/2\), so this gap is at most
\(\beta^{-1}\log2\).

More generally, under
\[
 \mu_n=(1-w_n)U_n+w_nU_{C_n},
\]
the same set sector has weight
\[
 p_n=w_n+(1-w_n)\delta_n.
\]
Assuming uniform-input AA, this family has full AA exactly when
\[
 \exists a,C>0\quad
 m\le Cw_n^{-a}\quad\text{at every index with }w_n>0.
\]
\end{proposition}

\begin{proof}
On the code \(P(c_A)=h(A)\), so
\[
 \Var_{U_C}P=r(1-r)=\frac{1-1/n}{4}\ge\frac3{16}.
\]
Lemma~\ref{ian:hard-sector-orbit-covariance}, with \(\eta\equiv1\)
and \(k=1\), supplies the covariance bound. Every Gibbs codeword has
the same weight, hence conditioning on \(C_n\) gives \(U_{C_n}\).
Summing the two energy levels proves the mass and free-energy formulas.
The exact mixture identity and \(w_{1,n}=1/2\) are
Section~\ref{qm:exact-mixture-identity-and-transition-width}.

For arbitrary mixture weights, conditioning on the code again gives
\(U_C\), with the displayed actual mass. The full-AA equivalence is
Theorem~\ref{threshold:weight-criterion} in
\path{obstructions/noisy_codes/04_energy_entropy_and_schedules.tex};
its sufficiency uses uniform-input AA, while its necessity uses the
one-query lower bound.
\end{proof}

A mixture alone has no specified temperature or free energy. If
\(0\le w_n<1\) is realized at a specified \(\beta>0\) by a two-level
energy, one may take
\[
 \Delta_n=\frac1\beta
 \log\left(1+\frac{w_n}{(1-w_n)\delta_n}\right).
\]
Its code-sector gap is exactly \(\beta^{-1}\log(1/p_n)\).
The case \(w_n=1\) is a limiting code-supported law. In particular,
the mixture weight must not be substituted for the actual set mass
in the free-energy identity.

\begin{proposition}[The nonmonotone sector]
\label{ian:hard-sector-nonmonotone-instance}
For the Hamiltonian in
Section~\ref{sec:resolution-nonmonotone}, let
\[
 H=C_n\setminus\{c_0\}.
\]
For \(n\ge16\) and \(t\le m/2-2\),
\[
 \operatorname{Hard}\left(
 P,\mu_{\beta,n}(\cdot\mid H);
 t,\frac29,\frac5{24}\right)
 \qquad(\beta>0).
\]
Its exact mass and gap are
\[
 \begin{split}
 Z_n&=1+(n-1)n^{-\beta/2}+(2^n-n)2^{-\beta n},\\
 p_H&=\frac{(n-1)n^{-\beta/2}}{Z_n},\\
 F_H-F
 &=\frac1\beta\log\left(
 1+\frac{1+(2^n-n)2^{-\beta n}}
          {(n-1)n^{-\beta/2}}\right).
 \end{split}
\]
\end{proposition}

\begin{proof}
The conditional law on \(H\) is uniform. Its acceptance mean is
\[
 r_H=\frac{nr}{n-1}
 =\frac{\sqrt n}{2(\sqrt n+1)},
\]
so \(r_H(1-r_H)\ge2/9\).

For a depth-\(t\) test \(F\), write \(F_0=F(c_0)\). Since
\(P(c_0)=0\), direct subtraction of this atom gives
\[
 \Cov_{U_H}(P,F)
 =\frac n{n-1}\Cov_{U_C}(P,F)
 +\frac{nr}{(n-1)^2}\bigl(F_0-\E_{U_C}F\bigr).
\]
Consequently,
\[
 |\Cov_{U_H}(P,F)|
 \le\frac n{n-1}c_{m,t}+\frac n{(n-1)^2}
 \le\frac{46}{225}<\frac5{24}<\frac29.
\]
The partition formulas follow by summing the three energy levels.
\end{proof}

The same repo proof also supplies the smooth sector \(U_C\), with
certified weight \(n^{1-\beta/2}/Z_n\), on the intermediate band.
The set sector above removes the exceptional ground-state atom.

For this fixed \(H\), at fixed \(\beta>0\),
\[
 F_H-F=
 \begin{cases}
 \displaystyle
 \frac{(1-\beta)n\log2-(1-\beta/2)\log n+o(1)}{\beta},
       &0<\beta<1,\\[2mm]
 o(1),&1\le\beta<2,\\
 \frac12\log2+o(1),&\beta=2,\\
 \displaystyle\frac{\beta/2-1}{\beta}\log n+o(1),
       &\beta>2.
 \end{cases}
\]
Thus \([1,2]\) is exactly the fixed-temperature band on which this
hard sector has a dimension-uniform free-energy gap. Its hidden signal
remains hard on the conditional law outside the band, but its mass
vanishes. Assuming uniform-input AA,
Section~\ref{sec:resolution-nonmonotone} proves full AA outside the
band and failure throughout the band, including both endpoints.
This is a hard free-energy band, with no monotonicity in inverse
temperature.

\begin{proposition}[Noisy and soft character codes]
\label{ian:hard-sector-soft-instance}
Let \(0<\rho<1\), and let \(\nu_{n,\rho}\) be the noisy signed-character
law of Theorem~\ref{threshold:noisy-query-theorem}. Whenever
\[
 \frac{2^t}{\sqrt n}\le\frac{\rho^4}{16},
\]
one has
\[
 \operatorname{Hard}\left(
 P,\nu_{n,\rho};
 t,\frac{3\rho^4}{16},\frac{\rho^4}{16}\right).
\]

Now fix \(\beta,L>0\), set
\[
 M=2n,\qquad \rho=\tanh(\beta L),\qquad
 S_L(x)=\sum_{v\in\{\pm c_a\}}e^{Lx\cdot v},
 \qquad \mu(x)=\frac{S_L(x)^\beta}{Z},
\]
and define
\[
 G=\left\{x:\exists v\ \forall w\ne v,\
 x\cdot v-x\cdot w\ge\rho n/2\right\}.
\]
Put
\[
 \begin{split}
 u_n&=(M-1)e^{-L\rho n/2},\\
 b_n&=(M-1)e^{-\rho^2n/16},\\
 e_n&=(M-1)\left[
 \beta2^{(\beta-1)_+}e^{-L\rho n/2}
 +e^{-\beta L\rho n/2}
 +M^{(\beta-1)_+}e^{-\rho^2n/16}\right],\\
 \tau_n&=\frac{e_n}{1-e_n}.
 \end{split}
\]
For every such \(n,t\) satisfying
\[
 u_n\le1,\qquad e_n<1,\qquad
 b_n+2\tau_n\le\rho^4/160,\qquad
 2^t/\sqrt n\le\rho^4/16,
\]
the set \(G\) satisfies
\[
 \operatorname{Hard}\left(
 P,\mu(\cdot\mid G);
 t,\frac{\rho^4}{8},\frac{\rho^4}{16}\right),
 \qquad \mu(G)\ge1-b_n-\tau_n.
\]
Its exact gap and a certified upper bound are
\[
 F_G-F
 =\frac1\beta\log
 \frac{\sum_xS_L(x)^\beta}{\sum_{x\in G}S_L(x)^\beta}
 \le\frac1\beta\log\frac1{1-b_n-\tau_n}.
\]
For each fixed \(\beta,L>0\), the displayed conditions hold for all
sufficiently large admissible \(n\), with a depth threshold
\(t=m/2-O_{\beta,L}(1)\).
\end{proposition}

\begin{proof}
In the noisy model, condition on the global sign and the entire
independent noise vector. Every visible query then uses one linear
query to uniform \(A\). Lemma~\ref{ian:hard-sector-orbit-covariance}
gives covariance at most \(\rho^4/32\), hence at most the stated
\(\rho^4/16\). The exact repo variance formula is
\[
 \Var_{\nu_{n,\rho}}P
 =r(1-r)\left[\rho^4+\frac{2(1-\rho^4)}n\right]
 \ge\frac{3\rho^4}{16}.
\]

For the soft model, the proof of
Theorem~\ref{threshold:soft-code-theorem} in
\path{obstructions/noisy_codes/03_fixed_hamiltonian_all_temperatures.tex}
defines
\[
 B(x)=\sum_ve^{\beta Lx\cdot v},\qquad
 R(x)=\frac{S_L(x)^\beta}{B(x)},\qquad
 q(x)=\frac{B(x)}{M[2\cosh(\beta L)]^n}.
\]
Here \(q=\nu_{n,\rho}\), and its explicit estimates give
\[
 q(G^c)\le b_n,\qquad
 \E_q|R-1|\le e_n,\qquad
 \|\mu-q\|_{\rm TV}\le\tau_n.
\]
The bound for \(e_n\) follows there from \(u_n\le1\). Hence
\[
 \mu(G^c)\le b_n+\tau_n,\qquad
 \|\mu(\cdot\mid G)-q\|_{\rm TV}\le b_n+2\tau_n.
\]

For laws \(\lambda,\lambda'\) at TV distance \(d\), a bounded
\(P\in[0,1]\) satisfies
\[
 |\Var_\lambda P-\Var_{\lambda'}P|\le d.
\]
Indeed, use
\(\Var_\lambda P=\min_{c\in[0,1]}\E_\lambda(P-c)^2\).
For \(F\in[-1,1]\), expansion of the two covariances gives
\[
 |\Cov_\lambda(P,F)-\Cov_{\lambda'}(P,F)|\le5d:
\]
the product expectation changes by at most \(2d\), the mean of \(P\)
by at most \(d\), and the mean of \(F\) by at most \(2d\).
Taking \(d=b_n+2\tau_n\) therefore gives
\[
 \Var_{\mu(\cdot\mid G)}P
 \ge\frac{3\rho^4}{16}-\frac{\rho^4}{160}
 >\frac{\rho^4}{8},
\]
and
\[
 |\Cov_{\mu(\cdot\mid G)}(P,F)|
 \le\frac{\rho^4}{32}+5\frac{\rho^4}{160}
 =\frac{\rho^4}{16}.
\]
The partition identity and gap estimate follow directly.
\end{proof}

There is also an explicit smooth-sector certificate, which does not
follow from TV comparison alone. The same repo proof gives the
pointwise bounds
\[
 r_-:=M^{\min\{\beta-1,0\}}\le R(x)
 \le M^{\max\{\beta-1,0\}},
 \qquad
 \mu(x)=q(x)\frac{R(x)}{\E_qR}.
\]
Thus
\[
 \mu\ge\frac{r_-}{\E_qR}\,q
 \ge\frac{r_-}{1+e_n}\,q.
\]
Whenever \(u_n\le1\), \(e_n<1\) and \(2^t/\sqrt n\le\rho^4/16\), the noisy law is consequently a smooth hard sector with certified
weight \(r_-/(1+e_n)\), and
\[
 \frac1\beta D_\infty(q\|\mu)
 \le\frac{\log(1+e_n)+(1-\beta)_+\log(2n)}{\beta}.
\]
At \(\beta=1\), \(\mu=q\) exactly. For \(0<\beta<1\) this certificate
has a polynomial loss in \(n\); the preceding set-sector argument
instead gives mass tending to one. The needed domination estimate
is therefore available from the repo's pointwise power bounds, even
though its stated comparison theorem is in TV distance.

\begin{proposition}[Local families, including every auxiliary and reference spin]
\label{ian:hard-sector-local-instances}
Let
\[
 G_C=\{x:\operatorname{dist}(x,C_n)<n/32\},\qquad
 G_\pm=\{x:\operatorname{dist}(x,\{\pm c_a\})<n/32\}.
\]
The following sectors have
\[
 \operatorname{Hard}\left(
 P,\mu_\beta(\cdot\mid H);
 t,\frac{17}{128},\frac18\right)
\]
for \(n\ge16\) and \(kt\le m/2-2\), whenever their displayed
\(\delta_n<1\). Their actual weights satisfy
\[
 p_H=\frac{Z_H}{Z}\ge1-\delta_n,\qquad
 0\le F_H-F\le\frac1\beta\log\frac1{1-\delta_n}.
\]

For the dense three-local family
\[
 E(X)=-\frac Jn\sum_{z,w}X_zX_wX_{z+w},
\]
take \(H=G_C\), \(k=1\), and
\[
 \delta_n=\frac1n
 \exp\bigl((\log2-\beta J/16)n\bigr).
\]
This tends to zero if \(\beta J>16\log2\).

For the sparse three-local family of
Theorem~\ref{res:sparse-theorem}, take \(H=G_C\), \(k=1\), and
\[
 \delta_n=\frac1n
 \exp\bigl((\log2-32\beta\kappa)n\bigr).
\]
This tends to zero if \(\beta\kappa>\log2/32\).

For the all-spin pairwise family of
Theorem~\ref{res:pairwise-theorem}, take
\[
 H=\{(X,S):X\in G_C\},\qquad k=3,
\]
and
\[
 \delta_n=\frac1n
 \exp\bigl((\log2-32K(\beta j))n\bigr),
 \qquad
 K(s)=\tfrac12\bigl(4\log\cosh(2s)-\log\cosh(4s)\bigr).
\]
This tends to zero if \(K(\beta j)>\log2/32\).

For the zero-field family of
Theorem~\ref{res:zero-field-theorem}, take
\[
 H=\{(X,S,R):X\in G_\pm\},\qquad k=3,
\]
and, with \(C=2048\),
\[
 \delta_n=\frac1{2n}
 \exp\bigl(((1+6C)\log2-64\beta j)n\bigr).
\]
This tends to zero if
\(\beta j>(1+6C)\log2/64\).
\end{proposition}

\begin{proof}
The mass estimates are the repo's respective energy-gap estimates in
\path{obstructions/local_codes/02_dense_three_local.tex},
\path{obstructions/local_codes/03_bounded_degree_sparsification.tex},
\path{obstructions/local_codes/04_pairwise_all_spins.tex}, and
\path{obstructions/local_codes/05_zero_field_pairwise.tex}.

On \(G_C\), the nearest character is unique. Character invariance
makes its label \(A\) uniform and independent of the residual noise
\(\eta=Xc_A\), conditional on the good event. The noise has fewer than
\(n/32\) negative coordinates.

In the pairwise model, conditional on \(X\), all auxiliary spins are
independent and each kernel depends on at most three visible spins,
by Lemma~\ref{res:gadget}. Realize these kernels using independent
uniform seeds. Fixing the seeds and \(\eta\) preserves uniformity of
\(A\). An auxiliary answer can then be computed using at most three
linear queries to \(A\).

In the zero-field model, the exact \((X,R)\) marginal is invariant
under \(X\mapsto Xc_b\) with \(R\) fixed. Conditional on the good
event, write uniquely
\[
 X=(-1)^Bc_A\eta.
\]
The signed character balls are disjoint, and \(A\) is uniform and
independent of \((B,\eta,R)\). Granting the entire reference vector
uses no queries to \(A\); fixing the independent auxiliary seeds
again makes every auxiliary answer computable with at most three
linear queries. These are exactly the all-spin reductions established
in the cited repo proofs.

Apply Lemma~\ref{ian:hard-sector-orbit-covariance} to each conditional
law. Its constants give dispersion \(17/128\), covariance at most
\(1/8\), and a positive difference \(1/128\). The mass and gap
statements follow from the partition identity.
\end{proof}

For clarity, the partition sums in this proposition are over all
physical spins. In the pairwise model, if \(E^{\rm pair}\) denotes
the sum of its thermal gadgets, they are
\[
 Z_H^{\rm pair}
 =\sum_{\substack{X\in G_C\\S}}e^{-\beta E^{\rm pair}(X,S)},
 \qquad
 Z^{\rm pair}=\sum_{X,S}e^{-\beta E^{\rm pair}(X,S)}.
\]
Writing \(M=|\mathcal E_n|\), the exact gadget identity gives
\[
 Z_H^{\rm pair}
 =e^{4Mc(\beta j)}
 \sum_{X\in G_C}
 \exp\left(K(\beta j)\sum_{(z,w)\in\mathcal E_n}
 X_zX_wX_{z+w}\right),
\]
with the same formula over all \(X\) for \(Z^{\rm pair}\).
The prefactor cancels in their ratio because the auxiliaries have
actually been summed.

For the zero-field energy \(E^{(0)}\) of the repo,
\[
 Z_H^{(0)}
 =\sum_{\substack{X\in G_\pm\\S,R}}
 e^{-\beta E^{(0)}(X,S,R)},\qquad
 Z^{(0)}=\sum_{X,S,R}e^{-\beta E^{(0)}(X,S,R)}.
\]
Equivalently,
\[
 Z_H^{(0)}
 =e^{4Mc(\beta j)}
 \sum_{\substack{X\in G_\pm\\R}}
 \exp\left(
 K(\beta j)\sum_{e=(z,w)}r_eX_zX_wX_{z+w}
 +\beta j\sum_{\{v,v'\}\in E(R_n)}r_vr_{v'}\right).
\]
Again the denominator uses every \(X,R\). No visible-spin partition
function may be substituted while discarding its auxiliary or
reference multiplicities.

All these sectors are hard in the variance--covariance sense of
Definition~\ref{ian:hard-sector-definition}. Their simulation
obstructions therefore follow from
Proposition~\ref{ian:hard-sector-forward}, and fit the asymptotic
characterizations of Theorem~\ref{ian:hard-sector-aa-equivalence} and Corollary~\ref{ian:hard-sector-resource-equivalence}.
